\documentclass[12pt,a4paper]{article}
\usepackage[utf8]{inputenc}
\usepackage[T1]{fontenc}
\usepackage[brazil]{babel}
\usepackage{geometry}
\usepackage{amsmath,amssymb}
\usepackage{booktabs}
\usepackage{longtable}
\usepackage{hyperref}
\usepackage{setspace}
\geometry{margin=2.5cm}
\onehalfspacing

\title{Adaptacao da Metodologia ITC de Potencial de Exportacoes\\para a Realidade Catarinense}
\author{Projeto Export Potential -- Santa Catarina}
\date{\today}

\begin{document}
\maketitle

\section{Objetivo do Relatorio}
Este relatorio documenta a adaptacao da metodologia do ITC (Export Potential Assessment) para o caso de Santa Catarina (SC), com foco em quatro blocos implementados no projeto:
\begin{itemize}
    \item \texttt{make\_demand.py} (componente de demanda),
    \item \texttt{make\_supply.py} (componente de oferta),
    \item \texttt{make\_ease.py} (componente de facilidade de comercio),
    \item \texttt{modeling/monetize\_epi\_sc.py} (monetizacao do indice e calculo de potencial realizado/nao realizado).
\end{itemize}

A logica central preserva a estrutura conceitual do ITC (demanda, oferta e facilidade), mas introduz adaptacoes para um recorte subnacional em que o horizonte de mercado mundial nao e usado como ancora final de valor monetario.

\section{Notacao}
A notacao adotada no relatorio e:
\begin{itemize}
    \item $i$: unidade exportadora (neste caso, $i=SC$),
    \item $j$: mercado importador,
    \item $k$: produto (SH6).
\end{itemize}

Variaveis principais:
\begin{itemize}
    \item $M_{jk}^{2024}$: importacao historica ponderada do mercado $j$ para o produto $k$,
    \item $D_j^{2030}$: indice de crescimento de demanda do mercado $j$ ate 2030,
    \item $\widehat{M}_{jk}^{2030}$: importacao projetada do mercado $j$ para o produto $k$,
    \item $X_{SC,k}^{w}$: exportacao historica ponderada de SC no produto $k$,
    \item $\widehat{X}_{SC,k}^{2030}$: exportacao projetada de SC no produto $k$,
    \item $s_{SC,k}^{2030}$: participacao projetada de SC no mercado mundial do produto $k$,
    \item $E_j$: facilidade de comercio de SC no mercado $j$,
    \item $B_{jk}$: exportacao bilateral observada (ponderada) de SC para $j$ no produto $k$.
\end{itemize}

\section{Referencial ITC e Adaptacoes}
\subsection{Estrutura original (ITC)}
A referencia ITC organiza o potencial de exportacoes em tres dimensoes:
\begin{enumerate}
    \item demanda potencial no destino,
    \item capacidade/competitividade de oferta do exportador,
    \item facilidade de acesso/comercio.
\end{enumerate}

\subsection{Adaptacoes para Santa Catarina}
Nesta adaptacao, os principais ajustes sao:
\begin{itemize}
    \item foco subnacional (SC) com base em participacao de SC nas exportacoes brasileiras por SH6,
    \item uso de projecoes de PIB e populacao para construir demanda ate 2030,
    \item calibracao de facilidade por historico bilateral revelado de SC,
    \item monetizacao final ancorada na oferta projetada de SC (e nao em um teto exogeno de mercado mundial agregado).
\end{itemize}

\section{Componente de Demanda (\texttt{make\_demand.py})}
\subsection{Passo 1: Projecao de populacao}
A partir de \texttt{pop\_growth.xlsx}, constroi-se o indice acumulado de populacao ate 2030:
\begin{equation}
P_j^{2030} = \prod_{t=2025}^{2030}(1+g_{j,t}^{pop})
\end{equation}

\subsection{Passo 2: Projecao de PIB}
Com \texttt{gdp\_growth.xlsx}, calcula-se:
\begin{equation}
Y_j^{2030} = \prod_{t=2025}^{2030}(1+g_{j,t}^{gdp})
\end{equation}

\subsection{Passo 3: PIB per capita e elasticidade-renda}
\begin{equation}
\left(\frac{Y}{P}\right)_j^{2030} = \frac{Y_j^{2030}}{P_j^{2030}}
\end{equation}

Aplica-se a elasticidade media da demanda em relacao ao PIB per capita ($\varepsilon=1{,}201$):
\begin{equation}
A_j^{2030} = \left(\left(\frac{Y}{P}\right)_j^{2030}\right)^{\varepsilon}
\end{equation}

Indice de demanda:
\begin{equation}
D_j^{2030} = A_j^{2030} \cdot P_j^{2030}
\end{equation}

\subsection{Passo 4: Importacao projetada por mercado-produto}
Com base no historico ponderado de importacoes ($M_{jk}^{2024}$):
\begin{equation}
\widehat{M}_{jk}^{2030} = M_{jk}^{2024} \cdot D_j^{2030}
\end{equation}

\section{Componente de Oferta (\texttt{make\_supply.py})}
\subsection{Passo 1: Reconstrucao da serie de SC}
Partindo das exportacoes brasileiras por SH6, usa-se a participacao de SC por produto e ano (arquivo \texttt{share\_sc.xlsx}):
\begin{equation}
X_{SC,jk,t} = X_{BRA,jk,t} \cdot share_{SC,k,t}
\end{equation}

Em seguida, calcula-se media ponderada de 8 anos para estabilizar o sinal de oferta:
\begin{equation}
X_{SC,k}^{w} = \frac{\sum_t \omega_t X_{SC,k,t}}{\sum_t \omega_t}
\end{equation}

\subsection{Passo 2: Projecao de oferta de SC para 2030}
No script, usa-se um fator acumulado de crescimento do PIB para SC:
\begin{equation}
\widehat{X}_{SC,k}^{2030} = X_{SC,k}^{w} \cdot 1{,}202
\end{equation}

\subsection{Passo 3: Participacao projetada de SC por SH6}
Tambem se projeta a oferta de todos os exportadores para 2030 e calcula-se a participacao relativa de SC:
\begin{equation}
s_{SC,k}^{2030} = \frac{\widehat{X}_{SC,k}^{2030}}{\sum_e \widehat{X}_{e,k}^{2030}}
\end{equation}

\section{Componente de Facilidade de Comercio (\texttt{make\_ease.py})}
\subsection{Passo 1: Fluxo bilateral revelado de SC}
Gera-se $B_{jk}$ como exportacao bilateral historica ponderada de SC por mercado-produto, e $B_j=\sum_k B_{jk}$.

\subsection{Passo 2: Benchmark potencial por destino}
Define-se um denominador de referencia por mercado:
\begin{equation}
S_j = \sum_k \left(M_{jk}^{2024} \cdot s_{SC,k}^{2030}\right)
\end{equation}

\subsection{Passo 3: Facilidade de comercio}
\begin{equation}
E_j = \frac{B_j}{S_j}
\end{equation}

Interpretacao: quanto maior $E_j$, maior a evidencia de que SC consegue converter potencial tecnico em comercio efetivo naquele destino.

\section{Indice de Potencial (base para monetizacao)}
No modulo de monetizacao, o score bruto por mercado-produto e:
\begin{equation}
W_{jk} = s_{SC,k}^{2030} \cdot \widehat{M}_{jk}^{2030} \cdot E_j
\end{equation}

Esse termo ja esta em unidade monetaria (US$ ajustados), pois combina uma base monetaria ($\widehat{M}_{jk}^{2030}$) com dois fatores adimensionais.

\section{Monetizacao do EPI e Potencial Nao Realizado (\texttt{monetize\_epi\_sc.py})}
\subsection{Passo 1: Share de alocacao dentro de cada SH6}
\begin{equation}
\alpha_{jk} = \frac{W_{jk}}{\sum_j W_{jk}}
\end{equation}

\subsection{Passo 2: Alocacao da oferta projetada de SC}
\begin{equation}
A_{jk} = \widehat{X}_{SC,k}^{2030} \cdot \alpha_{jk}
\end{equation}

\subsection{Passo 3: Teto de viabilidade de mercado}
\begin{equation}
V_{jk} = W_{jk}
\end{equation}

\subsection{Passo 4: Potencial monetario factivel}
\begin{equation}
P_{jk} = \min(A_{jk}, V_{jk})
\end{equation}

\subsection{Passo 5: Realizado, nao realizado e utilizacao}
Com $B_{jk}$ como exportacao observada:
\begin{align}
U_{jk} &= \max(P_{jk} - B_{jk}, 0) \\
R_{jk} &= \min(P_{jk}, B_{jk}) \\
O_{jk} &= \max(B_{jk} - P_{jk}, 0) \\
\rho_{jk} &= \begin{cases}
\frac{B_{jk}}{P_{jk}}, & P_{jk}>0 \\
0, & P_{jk}=0
\end{cases}
\end{align}

onde:
\begin{itemize}
    \item $U_{jk}$: potencial nao realizado,
    \item $R_{jk}$: potencial realizado,
    \item $O_{jk}$: comercio acima do potencial modelado,
    \item $\rho_{jk}$: taxa de utilizacao do potencial.
\end{itemize}

\section{Classificacao em Faixas para Visualizacao}
O script aplica \textit{k-means} (5 grupos) separadamente para \texttt{potential\_value} e \texttt{unrealized\_potential\_value}, com rotulos ordinais:
\begin{center}
Baixo, Medio, Medio-Alto, Alto, Muito Alto.
\end{center}

Essa etapa e de comunicacao/segmentacao e nao altera os valores monetarios estruturais do potencial.

\section{Hipoteses Metodologicas Principais}
\begin{enumerate}
    \item A demanda futura e bem representada por crescimento de populacao e PIB per capita com elasticidade fixa ($\varepsilon=1{,}201$).
    \item A competitividade relativa de SC por SH6 e capturada pela participacao projetada $s_{SC,k}^{2030}$.
    \item A facilidade de comercio e inferida por desempenho bilateral historico (efeito revelado).
    \item O potencial monetario final deve respeitar simultaneamente restricao de oferta (SC) e restricao de viabilidade de mercado.
    \item O modelo e comparativo-estatico e nao resolve equilibrio geral (preco relativo, substituicao global, feedback competitivo).
\end{enumerate}

\section{Intuicao Economica dos Componentes}
\subsection{Demanda}
A demanda indica \textit{onde o mercado cresce}. Se um pais tende a aumentar importacoes do produto, esse destino ganha peso no potencial.

\subsection{Oferta}
A oferta indica \textit{o que SC consegue entregar}. Mesmo com demanda alta, o potencial nao pode exceder a capacidade projetada de exportacao de SC por SH6.

\subsection{Facilidade de comercio}
A facilidade indica \textit{quao convertivel e a oportunidade}. Mercados com historico melhor de penetracao de SC recebem multiplicador maior.

\subsection{Intuicao geral do indicador}
O potencial monetario $P_{jk}$ e o valor factivel no par mercado-produto. Ele nao e previsao automatica de exportacao, nem incremento garantido.

\begin{itemize}
    \item \textbf{Potencial realizado} ($R_{jk}$): parte do potencial que ja esta sendo capturada.
    \item \textbf{Potencial nao realizado} ($U_{jk}$): espaco adicional, dentro das restricoes do modelo.
\end{itemize}

No agregado, a leitura gerencial recomendada e:
\begin{enumerate}
    \item tamanho da fronteira factivel ($\sum_{jk}P_{jk}$),
    \item quanto ja foi capturado ($\sum_{jk}R_{jk}$),
    \item quanto ainda e oportunidade ($\sum_{jk}U_{jk}$),
    \item onde ha risco de sustentacao ($\sum_{jk}O_{jk}$).
\end{enumerate}

\section{Produtos Gerados no Pipeline}
No estado atual do projeto, os principais produtos de saida da monetizacao sao:
\begin{itemize}
    \item \texttt{data/processed/epi\_monetary\_sc.parquet}
    \item \texttt{data/processed/epi\_monetary\_sc\_country.parquet}
    \item \texttt{data/processed/epi\_monetary\_sc\_sh6.parquet}
\end{itemize}

\end{document}
